\subsection{Kiểm định giả thuyết}

Ta có mệnh đề sau:

\begin{prop}
  Mọi level-$\alpha$ sequential test $\phi$ cho $\mathcal{P}$ đều có thể nhận
  được bằng cách lấy ngưỡng $1/\alpha$ một e-process $M$ nào đó cho
  $\mathcal{P}$.
\end{prop}

Do đó, các e-process cung cấp một khung lý thuyết hoàn chỉnh cho kiểm định
tuần tự. Thay vì tập trung vào chính kiểm định đó, chúng ta có thể tập trung
vào việc xây dựng các e-process. Tuy nhiên, chúng ta không chỉ coi e-processes
là một công cụ để suy ra các kiểm định. Chúng ta coi e-process là các thước đo
bằng chứng chống lại giả thuyết không: giá trị của chúng càng lớn (hoặc tăng
lên), chúng ta càng có nhiều bằng chứng chống lại giả thuyết không. Chúng ta
có thể đặt ngưỡng cho chúng để đưa ra quyết định nếu muốn, nhưng chúng vẫn giữ
được khả năng diễn giải ngay cả khi không cần đặt ngưỡng.

\subsubsection{Kiểm định qua cá cược}

Có một cách đơn giản để xem các test (super)martingale như các tiến trình tài
sản phát sinh từ các trò chơi. Ý tưởng chủ chốt như sau:

Để kiểm định một giả thuyết không $\mathcal{P}$, chúng ta thiết lập một trò
chơi may rủi, trong đó nhà thống kê  thực hiện đặt cược vào mỗi quan sát trước
khi quan sát nó. Trò chơi này phải thỏa mãn hai tính chất:
\begin{enumerate}[label=(\alph*)]
  \item Nếu giả thuyết không là đúng, nghĩa là $P \in \mathcal{P}$, thì không
    một chiến lược cá cược nào có thể kiếm tiền một cách hệ thống (tức là, nó
    có thể kiếm tiền do may mắn, nhưng không thể đảm bảo kiếm được tiền);

  \item Nếu giả thuyết không là sai, thì phải tồn tại các chiến lược cá cược
    ``tốt'' có thể kiếm được tiền (mặc dù người ta có thể không bao giờ chắc
    chắn về lợi nhuận trong ngắn hạn, nhưng có một cách để đảm bảo kiếm được
    tiền trong dài hạn).
\end{enumerate} Nếu một trò chơi như vậy có thể được thiết kế, thì tài sản của
nhà thống kê trong trò chơi đó trực tiếp là một thước đo bằng chứng chống lại
giả thuyết không: tài sản càng nhiều (hoặc chính xác hơn là tỉ số giữa tài sản
cuối cùng và tài sản ban đầu càng lớn), thì càng có nhiều bằng chứng cho thấy
giả thuyết không có khả năng là sai.

Chúng ta hãy bắt đầu với ví dụ đơn giản nhất, như đã thảo luận trong phần
trước. Chúng ta chỉ ra cách diễn giải việc kiểm định $P$ trong
\algorithmname~\ref{alg:simple-P}.

\begin{algorithm}[htbp]
  \caption{Kiểm định qua cá cược: trường hợp $P$ đơn}
  \label{alg:simple-P}
  Khởi tạo tài sản ban đầu cho nhà thống kê $W_0 = 1$\;
  \For{$t = 1$, $2$, $\ldots\,$}{
    Nhà thống kê tuyên bố một khoản cược $S_t \colon \mathcal{X} \to [0,
    \infty]$ sao cho \begin{equation*}
      \expect_P[S_t(X) \mid X_1, \dots, X_{t - 1}] \le 1
    \end{equation*}\vspace{-1\baselineskip}\;
    Quan sát $X_t$\;
    Cập nhật tài sản của nhà thống kê $W_t = W_{t - 1} \cdot S_t(X_t)$\;
  }
\end{algorithm}

Có ba điểm đáng chú ý cần bình luận ngay. Thứ nhất, trong trò chơi trên, ràng
buộc đối với $S_t$ đơn giản là nó là một e-variable cho $\mathcal{P}$ có điều
kiện dựa trên quá khứ. Điều này hàm ý ngay lập tức rằng đối với bất kỳ chiến
lược cá cược nào, tiến trình tài sản $W$ là một test supermartingale cho
$\mathcal{P}$.

Thứ hai, người ta có thể hỏi: đâu là khoản cược tối ưu? Giả sử để đơn giản rằng
giả thuyết thay thế $Q$ là đơn giản, và $P \ll Q$. Khi đó, tài sản có thể tăng
trưởng theo hàm mũ dưới $Q$, và do đó việc tối ưu hóa số mũ là tự nhiên. Điều
này có nghĩa là chọn $S_t$ để tối ưu hóa $\expect_Q[\log S_t(X) | X_1, \dots,
X_{t-1}]$ với điều kiện $S_t$ là một e-variable cho $\mathcal{P}$. Câu trả lời
đã được rút ra, đó là tỉ số khả năng (likelihood ratio) của $Q$ so với $P$, có
điều kiện dựa trên $X_1, \dots, X_{t-1}$ (phép điều kiện này không liên quan
nếu $P$ và $Q$ là các phân phối tích i.i.d trên $\mathcal{X}^\infty$, nhưng nó
lại quan trọng nếu chúng không phải như vậy).

Điểm cuối cùng là nếu người ta muốn kiểm định một giả thuyết không hợp
$\mathcal{P}$, thay đổi duy nhất đối với giao thức là yêu cầu $S_t$ phải là một
e-variable cho $\mathcal{P}$, nghĩa là ràng buộc trong giao thức trên phải
giữ nguyên với mọi $P \in \mathcal{P}$. Khi đó, tiến trình tài sản $W$ là một
testsupermartingale cho $\mathcal{P}$.

Bây giờ, chúng ta có thể tự nhiên thắc mắc làm thế nào các e-process có thể
được diễn giải trong khuôn khổ kiểm định qua cá cược này, trái ngược với việc
chỉ đơn thuần là các test supermartingale. Câu trả lời là diễn giải lý thuyết
trò chơi của chúng phức tạp hơn một chút. Nhà thống kê không chơi một trò chơi
duy nhất chống lại $\mathcal{P}$, mà thay vào đó chơi một trò chơi riêng biệt
chống lại mỗi $P \in \mathcal{P}$, và chấp nhận tài sản ròng của mình là tài
sản tệ nhất trong tất cả các trò chơi này. Điều này được làm rõ trong
\algorithmname\ref{alg:composite-P}.

\begin{algorithm}[htbp]
  \caption{Kiểm định qua cá cược: trường hợp $P$ hợp}
  \label{alg:composite-P}

  \For{$P \in \mathcal{P}$}{
    Khởi tạo tài sản ban đầu của nhà thống kê cho trò chơi $P$ $W_0^P = 1$\;
  }
  \BlankLine
  \For{$t = 1$, $2$, $\ldots\,$}{
    Nhà thống kê tuyên bố một khoản cược $S^P_t \colon \mathcal{X} \to [0,
    \infty]$ sao cho \begin{equation*}
      \expect_P[S^P_t(X) \mid X_1, \dots, X_{t - 1}] \le 1
    \end{equation*}\vspace{-1\baselineskip}\;
    Quan sát $X_t$\;
    Cập nhật tài sản trong $P$ của nhà thống kê $W^P_t = W^P_{t - 1} \cdot
    S^P_t(X_t)$\;

    Tài sản tổng thể của nhà thống kê được cập nhật $W_t = \inf_{P \in
    \mathcal{P}} W^P_t$\;
  }
\end{algorithm}

\subsubsection{Trường hợp giả thuyết không và thay thế đều đơn}

Giả sử chúng ta đang kiểm định một giả thuyết đơn $P$ đối với một giả thuyết
đơn $Q \ll P$. Nếu chúng ta quan sát dữ liệu i.i.d $X_1, X_2, \dots$ một cách
tuần tự, thì tiến trình tỉ số khả năng $M$ được cho bởi:
\begin{equation}
    M_0 = 1 \text{ và } M_t = \prod_{i=1}^{t} \frac{dQ}{dP}(X_i) \text{ với } t
    \in \mathbb{N},
\end{equation}
là một e-process thích nghi với sự lọc được tạo ra bởi dữ liệu, trong đó
$dQ/dP(X)$ là tỉ số khả năng của $Q$ đối với $P$ khi quan sát $X$. Chúng ta đã
bắt gặp e-process này trong Mục 7.2, được sử dụng trong SPRT. Hơn nữa, dễ dàng
thấy rằng $M$ là một test martingale cho $P$. Trên thực tế, ngay cả khi dữ liệu
không phải i.i.d, tỉ số khả năng dựa trên $n$ điểm dữ liệu đầu tiên:
\begin{equation}
    M_n = \frac{dQ}{dP}(X_1, \dots, X_n),
\end{equation}
vẫn là một test martingale cho $P$ (có thể kiểm tra bằng phép lấy điều kiện).

Khi kiểm định một giả thuyết không đơn giản $P$, việc sử dụng một test
martingale để xây dựng các e-process là tối ưu: Bất kỳ e-process nào cho $P$
đều có thể bị chặn trên bởi Snell envelope của nó dưới $P$ (vốn là một test
supermartingale), và đến lượt nó, lại có thể bị chi phối bởi test martingale
cho $P$ được cho bởi phân tích Doob.

Trong bối cảnh kiểm định giả thuyết không đơn $P$ đối với $Q \ll P$, tỉ số
khả năng $M_t$ là tối ưu logarit trong số tất cả các e-variable sử dụng $t$
điểm dữ liệu đầu tiên với mọi $t \in \mathbb{N}$.

Khi đó, khoản cược log-optimal là \begin{equation}
  S_t(X) = \frac{dQ}{dP}(X) = \frac{q(X)}{p(X)},
\end{equation} trong đó có đẳng thức sau nếu $P$, $Q$ các các hàm mật độ xác
suất $p$ và $q$ tương ứng. Khi đó $\expect_P[W_\tau] \le 1$, $\expect_P[\log
W_\tau] \le 0$. Trong khi đó, $\expect_Q[\log W_T] > 0$ được lấy tối đa và
tương đương $\KL(Q, P)$.

\subsubsection{Các trường hợp khác}

Các trường hợp giả thuyết không hợp, giả thuyết thay thế đơn; giả thuyết
không đơn, giả thuyết thay thế hợp; hay cả giả thuyết không và thay thế đều hợp
đều khá phức tạp. Do đó, báo cáo này không trình bày các phần này. Ta có thể
tham khảo thêm trong \cite{ramdas2023savi}
hoặc \cite{ramdas2025ebook}.
